\documentclass[letterpaper, 10pt, conference]{ieeeconf}
\IEEEoverridecommandlockouts
\overrideIEEEmargins

\usepackage{amsmath,amssymb}
\usepackage{graphicx}
\usepackage{booktabs}
\usepackage{siunitx}
\usepackage{xurl}
\usepackage[hidelinks]{hyperref}
\usepackage{placeins}

\hyphenpenalty=10000
\exhyphenpenalty=10000
\emergencystretch=2em

% shorthands
\newcommand{\q}{\mathbf{q}}
\newcommand{\w}{\boldsymbol{\omega}}
\newcommand{\ez}{\mathbf{e}_z}
\newcommand{\Hstar}{\mathbb{H}_*}
\newcommand{\dg}{^\circ}
\DeclareMathOperator{\sinc}{sinc}

\newcommand{\figorbox}[2][\linewidth]{%
  \includegraphics[width=#1]{#2}%
}

\begin{document}

\title{ECE 276A Project 1: Orientation Tracking}

\author{\authorblockN{Aditya Wirayutha Kaya Simto}
\authorblockA{Department of Computer Science and Systems Engineering \\ Kobe University}}

\maketitle
\thispagestyle{empty}
\pagestyle{empty}

\section{Introduction}

This project estimates the orientation of an IMU over time. The gyroscope gives
angular velocity, which can be integrated into orientation, but sensor bias
causes drift. The accelerometer gives the gravity direction when linear motion
is small. Gravity can correct roll and pitch, but it cannot correct yaw. The
main problem is therefore to combine both sensors without losing the useful
information from either one.

The data contains eleven sequences. Nine include VICON ground truth and two are
test sequences. Following the handout~\cite{atanasov}, I calibrate the raw IMU
data, integrate the gyroscope to obtain an initial quaternion trajectory, and
refine it with projected gradient descent. I then use the estimated orientations
to project the camera images into an equirectangular panorama.

\section{Problem Formulation}

At time $t$, $\q_t\in\Hstar$ describes the body orientation relative to the
world. My quaternion functions store $\q=(q_w,\mathbf{q}_v)$ as
$[q_w,q_x,q_y,q_z]$. Before passing a quaternion to SciPy in
\texttt{main.py}, I reorder it as $[q_x,q_y,q_z,q_w]$ because SciPy uses the
scalar last order by default~\cite{scipyquat}. The IMU gives angular velocity
$\w_t$ and acceleration $\mathbf a_t$ in the body frame. Since the sampling
interval changes slightly, I use the measured interval $\tau_t$ at every step.

\emph{Motion model.} I assume $\w_t$ is constant during one interval:
\begin{equation}
\q_{t+1}=f(\q_t,\tau_t\w_t):=\q_t\circ\exp\!\left(\left[0,\tfrac{\tau_t\w_t}{2}\right]\right),
\label{eq:motion}
\end{equation}
where $\circ$ is quaternion multiplication. The increment is on the right
because $\w_t$ is measured in the body frame.

\emph{Observation model.} When linear acceleration is small, the accelerometer
measures gravity in the body frame:
\begin{equation}
[0,\mathbf{a}_t]=h(\q_t):=\q_t^{-1}\circ[0,0,0,g]\circ\q_t,
\label{eq:obs}
\end{equation}
where $g=1$ in gravity units. The same model is
$h(\q)=R(\q)^\top\ez$, with $\ez=[0,0,1]^\top$.

\emph{Cost.} I estimate the trajectory by minimizing both model errors:
\begin{equation}
\begin{aligned}
\mathbf r_t^m&=2\log\!\left(
\q_{t+1}^{-1}\circ f(\q_t,\tau_t\w_t)\right),\\
\mathbf r_t^a&=[0,\mathbf a_t]-h(\q_t),\\
c(\q_{1:T})&=\tfrac12\sum_{t=0}^{T-1}\|\mathbf r_t^m\|_2^2
+\tfrac12\sum_{t=0}^{T}\|\mathbf r_t^a\|_2^2.
\end{aligned}
\label{eq:cost}
\end{equation}
The first term follows the gyroscope motion. The second term follows gravity.
The code includes $t=0$ in the gravity term. Since $\q_0$ is fixed, this adds
the same initial error to every candidate trajectory.

\emph{Constraint.} Each quaternion must have unit norm:
\begin{equation}
\min_{\q_{1:T}}c(\q_{1:T})\quad\text{s.t.}\quad\|\q_t\|_2=1\ \forall t.
\label{eq:constrained}
\end{equation}
The constraint set is the unit sphere $S^3\subset\mathbb{R}^4$. Its Euclidean
projection is
\begin{equation}
\Pi_{\Hstar}(\q)=\q/\|\q\|_2.
\label{eq:proj}
\end{equation}
This projection is only quaternion normalization. I fix the initial orientation
as $\q_0=[1,0,0,0]$. Quaternions are convenient here because rotation matrices
need more constraints and Euler angles become singular at pitch $\pm90\dg$.

\section{Technical Approach}

\subsection{Orientation Tracking}

\subsubsection{Data Preprocessing}

The IMU file is a $7\times N$ array containing time, acceleration, and
gyroscope data. In \texttt{find\_static\_end}, I call the NumPy
\texttt{sliding\_window\_view} function~\cite{numpywindow} with 50 samples to
find when the sensor starts moving:
\begin{equation}
n_s=\min\left\{k:
\max_{d\in\{x,y,z\}}\operatorname{std}(G_{k:k+49,d})>1.5\right\}.
\label{eq:static}
\end{equation}
The starting time is different in each dataset, so one fixed value would not
work well. Calling moving data stationary would make the calibration confidently
wrong, which is not the confidence I need here.

Here $S_a$ and $S_g$ are the sensor sensitivities, and
$V_{\mathrm{ref}}$ is the reference voltage. I compute the scales and biases as
\begin{equation}
\begin{aligned}
s_a&=\frac{V_{\mathrm{ref}}}{1023S_a}, &
s_g&=\frac{V_{\mathrm{ref}}}{1023S_g}\frac{\pi}{180},\\
\mathbf b_g&=\overline{\mathbf G}_{1:n_s}, &
\mathbf b_a&=\overline{\mathbf A}_{1:n_s}
 -\begin{bmatrix}0&0&1/s_a\end{bmatrix}^{\top}.
\end{aligned}
\label{eq:calibration}
\end{equation}
I then calibrate each measurement using
\begin{equation}
\mathbf a_t=s_a(\mathbf A_t-\mathbf b_a),\qquad
\w_t=s_g(\mathbf G_t-\mathbf b_g).
\label{eq:calibrated}
\end{equation}
This converts raw counts into g and rad/s. I estimate a separate bias for each
dataset.

The VICON and IMU timestamps do not match exactly. I use spherical linear
interpolation to place the VICON rotations at the IMU times, and I remove
invalid rotations. I check the calibration with
\begin{equation}
\frac{1}{|\mathcal I|}\sum_{t\in\mathcal I}
\left\|\mathbf a_t-R_{\mathrm{vicon}}(t)^\top\ez\right\|_2<0.1\ \text{g}.
\label{eq:gate}
\end{equation}
The residual is below the limit for every training dataset. The $+\ez$ sign
also gives much less error than $-\ez$, so I use $+\ez$. Thankfully, this sign
choice was not a close election.

\subsubsection{Optimization}

I first integrate the gyroscope. One rotation step and the resulting dead
reckoning trajectory are
\begin{equation}
\begin{aligned}
E_t&=\exp\!\left(\left[0,\frac{\tau_t\w_t}{2}\right]\right),\\
\q_0^{\mathrm{dr}}&=[1,0,0,0],\\
\q_{t+1}^{\mathrm{dr}}&=\q_t^{\mathrm{dr}}\circ E_t.
\end{aligned}
\label{eq:deadreckon}
\end{equation}
For a unit quaternion, $\|2\log\q\|=2\arccos(q_w)$. I define
$\mathbf f_t=\q_t\circ E_t$ and
$s_t=\langle\q_{t+1},\mathbf f_t\rangle$. This gives the simpler cost
\begin{equation}
c(\q)=2\sum_t\arccos^2(s_t)
+\frac12\sum_t\|\mathbf a_t-h(\q_t)\|_2^2.
\label{eq:cmotion}
\end{equation}
I clip $s_t$ to $[-1,1]$ before \texttt{arccos} to avoid numerical problems.
The mathematics already knows this range, but floating point arithmetic
occasionally needs a reminder.

For $\theta_t=\arccos(s_t)$, define
$\gamma_t=-4\theta_t/\sin\theta_t$. I add the motion gradient to both
quaternions in each time step:
\begin{equation}
\begin{aligned}
\mathbf g_{t+1}&\leftarrow
\mathbf g_{t+1}+\gamma_t\mathbf f_t,\\
\mathbf g_t&\leftarrow
\mathbf g_t+\gamma_t(\q_{t+1}\circ E_t^{-1}).
\end{aligned}
\label{eq:motiongrad}
\end{equation}
Direct division is unstable when $\theta_t=0$, so I use
$\gamma_t=-4/\sinc(\theta_t/\pi)$~\cite{numpy}. The gravity model and its
Jacobian are
\begin{equation}
\begin{aligned}
h(\q)&=\begin{bmatrix}
2(q_xq_z-q_wq_y)\\
2(q_yq_z+q_wq_x)\\
q_w^2-q_x^2-q_y^2+q_z^2
\end{bmatrix},\\
J(\q)&=2\begin{bmatrix}
-q_y&q_z&-q_w&q_x\\
q_x&q_w&q_z&q_y\\
q_w&-q_x&-q_y&q_z
\end{bmatrix},\\
\mathbf g_t&\leftarrow\mathbf g_t-J(\q_t)^\top
(\mathbf a_t-h(\q_t)).
\end{aligned}
\label{eq:obsgrad}
\end{equation}
Each optimization step moves against the gradient and normalizes the result:
\begin{equation}
\begin{aligned}
\q^{(k+1)}_t
&=\frac{\q^{(k)}_t-\alpha\mathbf g_t}
{\|\q^{(k)}_t-\alpha\mathbf g_t\|_2},\\
\alpha&=10^{-3},\qquad k=0,\ldots,1499.
\end{aligned}
\label{eq:pgd}
\end{equation}
I keep the first pose fixed by setting $\mathbf g_0=0$. I also compare the
analytic gradient with finite differences. The cost decreases smoothly and is
almost stable before the last iteration, as shown in Figure~\ref{fig:conv}.

\subsection{Panorama}

\subsubsection{Data Preprocessing}

The camera looks along body $x$. Body $y$ points left and body $z$ points up.
For pixel $(i,j)$, I define
\begin{equation}
\begin{aligned}
c_x&=\frac{W-1}{2}, & c_y&=\frac{H-1}{2},\\
u_j&=\frac{j-c_x}{W}, & v_i&=\frac{i-c_y}{H}.
\end{aligned}
\end{equation}
Image columns increase to the right and rows increase downward. These directions
are opposite to body $y$ and $z$, so both angles have a minus sign:
\begin{equation}
\lambda_{ij}=-u_j\operatorname{FOV}_h,\qquad
\varphi_{ij}=-v_i\operatorname{FOV}_v.
\end{equation}
Here $\lambda$ is the horizontal angle and $\varphi$ is the vertical angle.
I use $\operatorname{FOV}_h=60\dg$ and $\operatorname{FOV}_v=45\dg$, matching
the constants in \texttt{panorama.py}. The corresponding body ray is
\begin{equation}
\mathbf r_{ij}=
\begin{bmatrix}
\cos\varphi_{ij}\cos\lambda_{ij}\\
\cos\varphi_{ij}\sin\lambda_{ij}\\
\sin\varphi_{ij}
\end{bmatrix},\qquad \|\mathbf r_{ij}\|_2=1.
\label{eq:ray}
\end{equation}
At the image centre, $\lambda=\varphi=0$ and
$\mathbf r=[1,0,0]^\top$. For each camera frame, I use the latest IMU
orientation available at that time:
\begin{equation}
k(c)=\max\{k:t_k^{\mathrm{IMU}}\leq t_c^{\mathrm{cam}}\}.
\label{eq:timealign}
\end{equation}
This choice prevents the camera frame from using a future orientation.

\subsubsection{Projection and Computational Optimization}

I rotate each body ray into the world frame:
\begin{equation}
\mathbf p_{ijc}=R_{k(c)}\mathbf r_{ij}
=\begin{bmatrix}p_x&p_y&p_z\end{bmatrix}^{\top}.
\end{equation}
The world azimuth and elevation are
\begin{equation}
\theta=\operatorname{atan2}(p_y,p_x),\qquad
\psi=\arcsin(p_z),
\end{equation}
where $\theta\in(-\pi,\pi]$ and $\psi\in[-\pi/2,\pi/2]$. In the code, $p_z$ is
clipped to $[-1,1]$ before \texttt{arcsin}. For $H_p=480$ and $W_p=960$, I
convert these angles to panorama column $m$ and row $n$ with the equirectangular
projection~\cite{proj}:
\begin{equation}
\begin{aligned}
m&=\left\lfloor\frac{\theta+\pi}{2\pi}(W_p-1)\right\rfloor,\\
n&=\left\lfloor\frac{\pi/2-\psi}{\pi}(H_p-1)\right\rfloor.
\end{aligned}
\label{eq:equirectangular}
\end{equation}
The update is $P[n,m]\leftarrow I_c[i,j]$. When pixels overlap, I keep the most
recent one. There is no numerical optimizer in this part. I compute all camera
rays once, so each frame only needs one matrix multiplication. I ignore the
camera offset because the scene is far away, and I use an assumed field of view.
The code implements the equations directly with NumPy and does not call PROJ.
Figure~\ref{fig:smallsteps} shows four examples using values from dataset~1.

\begin{figure*}[t]
\centering
\begin{minipage}[t]{0.48\textwidth}
\centering
\figorbox{figures/small_static_window.png}
\end{minipage}
\hfill
\begin{minipage}[t]{0.48\textwidth}
\centering
\figorbox{figures/small_calibration.png}
\end{minipage}

\vspace{4pt}

\begin{minipage}[t]{0.48\textwidth}
\centering
\figorbox{figures/small_motion_step.png}
\end{minipage}
\hfill
\begin{minipage}[t]{0.48\textwidth}
\centering
\figorbox{figures/small_panorama_pixel.png}
\end{minipage}
\caption{Small numerical examples from dataset~1. The panels show static window
detection, IMU calibration, one dead reckoning update, and one camera pixel
projected into the panorama.}
\label{fig:smallsteps}
\end{figure*}

\section{Results}
\label{sec:results}

\subsection{Orientation Tracking}

Table~\ref{tab:acc} compares the estimates with VICON. After projected gradient
descent, the mean roll error decreases from $11.94\dg$ to $1.62\dg$, and the
mean pitch error decreases from $7.09\dg$ to $1.48\dg$. Roll improves in eight
datasets and pitch improves in all nine. Yaw improves by only $1.93\dg$ because
the accelerometer gives no yaw information. Figure~\ref{fig:peraxis} shows this
same pattern.

The final roll and pitch errors are between $0.98\dg$ and $2.18\dg$. Their
correlation with the accelerometer calibration residual is $r=0.903$, as shown
in Figure~\ref{fig:gate}. The complete orientation curves are included in
Appendix~\ref{app:perdataset}, Figures~\ref{fig:rpy} to~\ref{fig:rpy6to9}.

\begin{table*}[t]
\centering
\caption{Mean absolute error vs.\ VICON (degrees), gyroscope only $\to$ PGD.}
\label{tab:acc}
\begin{tabular}{lcccc}
\toprule
ds            & roll: gyro$\to$PGD        & pitch: gyro$\to$PGD       & yaw: gyro$\to$PGD          & geodesic: gyro$\to$PGD\\
\midrule
1             & 15.24 $\to$ \textbf{1.44} & 8.10 $\to$ \textbf{1.01}  & 21.04 $\to$ 13.19          & 17.08 $\to$ 12.98\\
2             & 14.99 $\to$ \textbf{1.28} & 14.41 $\to$ \textbf{1.08} & 18.97 $\to$ 13.32          & 22.65 $\to$ 13.50\\
3             & 1.26 $\to$ 1.77           & 2.28 $\to$ \textbf{1.86}  & 14.20 $\to$ 14.46          & 14.64 $\to$ 14.82\\
4             & 15.87 $\to$ \textbf{2.18} & 5.36 $\to$ \textbf{1.40}  & 38.90 $\to$ 39.04          & 42.25 $\to$ 39.16\\
5             & 18.57 $\to$ \textbf{1.90} & 9.03 $\to$ \textbf{2.11}  & 18.71 $\to$ 16.74          & 27.46 $\to$ 17.22\\
6             & 8.27 $\to$ \textbf{1.83}  & 3.07 $\to$ \textbf{1.52}  & 4.68 $\to$ 4.68            & 11.15 $\to$ 5.90\\
7             & 12.72 $\to$ \textbf{2.08} & 7.99 $\to$ \textbf{1.80}  & 6.28 $\to$ 5.02            & 15.57 $\to$ 5.83\\
8             & 3.19 $\to$ \textbf{1.10}  & 2.00 $\to$ \textbf{1.31}  & 10.47 $\to$ 10.11          & 11.34 $\to$ 10.49\\
9             & 17.38 $\to$ \textbf{0.98} & 11.61 $\to$ \textbf{1.25} & 29.61 $\to$ 28.89          & 36.46 $\to$ 29.09\\
\midrule
\textbf{mean} & \textbf{11.94 $\to$ 1.62} & \textbf{7.09 $\to$ 1.48}  & \textbf{18.10 $\to$ 16.16} &\\
\bottomrule
\end{tabular}
\end{table*}

\begin{figure*}[!t]
\centering
\figorbox{figures/summary_peraxis.png}
\caption{Mean roll, pitch, and yaw error before and after projected gradient
descent for all nine training datasets.}
\label{fig:peraxis}
\end{figure*}

\begin{figure*}[!t]
\centering
\figorbox{figures/summary_convergence.png}
\caption{Cost falls monotonically $131.6\to4.34$ (left); geodesic error vs.\
time (right).}
\label{fig:conv}
\end{figure*}

\begin{figure}[t]
\centering
\figorbox{figures/summary_gate_vs_error.png}
\caption{Accelerometer calibration residual and final roll and pitch error for
the nine training datasets.}
\label{fig:gate}
\end{figure}

Dataset~3 is the only case where the roll error increases, from $1.26\dg$ to
$1.77\dg$. Its gyroscope estimate is already good, so accelerometer noise moves
the result slightly away from it. Dataset~4 still has about $39\dg$ of yaw
error, which is expected because gravity cannot correct yaw.

I also tried fitting new sensor scales from VICON. The accelerometer residual
changed only from $0.0350$ to $0.0339$~g, while the gyroscope fit became
unstable. Using the fitted scales increased the errors of datasets 3 and 9 to
$60.06\dg$ and $69.25\dg$. For this reason, I kept the datasheet values. The
cost also decreases for test datasets 10 and 11, but their accuracy cannot be
measured because VICON is not available. Their curves are shown in
Figure~\ref{fig:rpytest}.

\subsection{Panorama}

The estimated orientations produce recognizable room panoramas.
Figure~\ref{fig:pano} compares VICON and PGD for dataset~1. The room layout is
similar, but the whiteboard appears twice in the PGD result. This matches the
remaining $13.19\dg$ yaw error. The room did not actually grow a second
whiteboard. Panorama coverage is between $31.6\%$ and
$80.4\%$. More training examples are shown in Figure~\ref{fig:panotrain}, and
the test results are shown in Figure~\ref{fig:panotest}.

\begin{figure*}[!t]
\centering
\figorbox{figures/summary_panorama.png}
\caption{VICON vs.\ estimated panorama; the doubled whiteboard is residual yaw.}
\label{fig:pano}
\end{figure*}

\begin{figure*}[!t]
\centering
\begin{minipage}[t]{0.49\textwidth}
\centering
\textbf{Dataset 2, VICON}\\[2pt]
\figorbox{figures/pano_ds2_vicon.png}
\end{minipage}
\hfill
\begin{minipage}[t]{0.49\textwidth}
\centering
\textbf{Dataset 2, PGD}\\[2pt]
\figorbox{figures/pano_ds2_pgd.png}
\end{minipage}

\begin{minipage}[t]{0.49\textwidth}
\centering
\textbf{Dataset 8, VICON}\\[2pt]
\figorbox{figures/pano_ds8_vicon.png}
\end{minipage}
\hfill
\begin{minipage}[t]{0.49\textwidth}
\centering
\textbf{Dataset 8, PGD}\\[2pt]
\figorbox{figures/pano_ds8_pgd.png}
\end{minipage}

\begin{minipage}[t]{0.49\textwidth}
\centering
\textbf{Dataset 9, VICON}\\[2pt]
\figorbox{figures/pano_ds9_vicon.png}
\end{minipage}
\hfill
\begin{minipage}[t]{0.49\textwidth}
\centering
\textbf{Dataset 9, PGD}\\[2pt]
\figorbox{figures/pano_ds9_pgd.png}
\end{minipage}
\caption{Panoramas from VICON and projected gradient descent for training
datasets 2, 8, and 9.}
\label{fig:panotrain}
\end{figure*}

\begin{figure*}[!t]
\centering
\textbf{Dataset 10}\\[2pt]
\figorbox[0.82\textwidth]{figures/pano_ds10_pgd.png}\\[6pt]
\textbf{Dataset 11}\\[2pt]
\figorbox[0.82\textwidth]{figures/pano_ds11_pgd.png}
\caption{Panoramas from projected gradient descent for test datasets 10 and
11. VICON is not available.}
\label{fig:panotest}
\end{figure*}

\section{Discussion}
\label{sec:discussion}

The main limitation is yaw. Rotating around gravity does not change the
accelerometer measurement, so yaw depends only on the gyroscope and continues
to drift. A magnetometer or visual reference would be needed to correct it.
The method also uses the full trajectory, including future measurements. This
helps the estimate, but it is not suitable for real time control.

Some settings are chosen manually. These include the camera field of view, the
static threshold, and equal weights for both cost terms. The model also assumes
that acceleration mainly comes from gravity. Fast linear motion breaks this
assumption. Dataset~3 also suggests that using sensor noise to choose the cost
weights may give a better result.

\FloatBarrier
\appendices
\section{Orientation Curves for Each Dataset}
\label{app:perdataset}

This appendix shows the roll, pitch, and yaw curves for every dataset. The main
text only uses the summary plots.

\begin{figure*}[!t]
\centering
\figorbox[0.88\textwidth]{figures/rpy_ds1.png}
\caption{Roll, pitch, and yaw for dataset~1. The plot compares VICON, gyroscope
integration, and projected gradient descent.}
\label{fig:rpy}
\end{figure*}

\begin{figure*}[!t]
\centering
\begin{minipage}[t]{0.49\textwidth}
\centering
\textbf{Dataset 2}\\[2pt]
\figorbox[\linewidth]{figures/rpy_ds2.png}
\end{minipage}
\hfill
\begin{minipage}[t]{0.49\textwidth}
\centering
\textbf{Dataset 3}\\[2pt]
\figorbox[\linewidth]{figures/rpy_ds3.png}
\end{minipage}

\begin{minipage}[t]{0.49\textwidth}
\centering
\textbf{Dataset 4}\\[2pt]
\figorbox[\linewidth]{figures/rpy_ds4.png}
\end{minipage}
\hfill
\begin{minipage}[t]{0.49\textwidth}
\centering
\textbf{Dataset 5}\\[2pt]
\figorbox[\linewidth]{figures/rpy_ds5.png}
\end{minipage}
\caption{Roll, pitch, and yaw for training datasets 2 to 5. Each plot compares
VICON, gyroscope integration, and projected gradient descent.}
\label{fig:rpy2to5}
\end{figure*}

\begin{figure*}[!t]
\centering
\begin{minipage}[t]{0.49\textwidth}
\centering
\textbf{Dataset 6}\\[2pt]
\figorbox[\linewidth]{figures/rpy_ds6.png}
\end{minipage}
\hfill
\begin{minipage}[t]{0.49\textwidth}
\centering
\textbf{Dataset 7}\\[2pt]
\figorbox[\linewidth]{figures/rpy_ds7.png}
\end{minipage}

\begin{minipage}[t]{0.49\textwidth}
\centering
\textbf{Dataset 8}\\[2pt]
\figorbox[\linewidth]{figures/rpy_ds8.png}
\end{minipage}
\hfill
\begin{minipage}[t]{0.49\textwidth}
\centering
\textbf{Dataset 9}\\[2pt]
\figorbox[\linewidth]{figures/rpy_ds9.png}
\end{minipage}
\caption{Roll, pitch, and yaw for training datasets 6 to 9. Each plot compares
VICON, gyroscope integration, and projected gradient descent.}
\label{fig:rpy6to9}
\end{figure*}

\begin{figure*}[!t]
\centering
\begin{minipage}[t]{0.49\textwidth}
\centering
\textbf{Dataset 10}\\[2pt]
\figorbox[\linewidth]{figures/rpy_ds10.png}
\end{minipage}
\hfill
\begin{minipage}[t]{0.49\textwidth}
\centering
\textbf{Dataset 11}\\[2pt]
\figorbox[\linewidth]{figures/rpy_ds11.png}
\end{minipage}
\caption{Roll, pitch, and yaw from gyroscope integration and projected gradient
descent for test datasets 10 and 11. VICON is not available.}
\label{fig:rpytest}
\end{figure*}

\clearpage

\begin{thebibliography}{5}
\bibitem{atanasov} N.~Atanasov, ``ECE 276A: Sensing \& Estimation in
Robotics,'' University of California San Diego. [Online]. Available:
\url{https://natanaso.github.io/ece276a/index.html}
\bibitem{numpy} NumPy Developers, ``\texttt{numpy.sinc}.'' [Online]. Available:
\url{https://numpy.org/doc/stable/reference/generated/numpy.sinc.html}
\bibitem{numpywindow} NumPy Developers,
``\texttt{numpy.lib.stride\_tricks.sliding\_window\_view}.'' [Online]. Available:
\url{https://numpy.org/doc/stable/reference/generated/numpy.lib.stride_tricks.sliding_window_view.html}
\bibitem{scipyquat} SciPy Developers, ``\texttt{Rotation.as\_quat}.'' [Online]. Available:
\url{https://docs.scipy.org/doc/scipy/reference/generated/scipy.spatial.transform.Rotation.as_quat.html}
\bibitem{proj} PROJ Contributors, ``Equidistant Cylindrical Projection.'' [Online]. Available:
\url{https://proj.org/en/stable/operations/projections/eqc.html}
\end{thebibliography}

\end{document}
